% ============================================================
% CHAPTER 5: k-CENTER (matches Vazirani Ch. 5)
% ============================================================
\setcounter{chapter}{4}
\chapter{k-Center}
\label{ch:kcenter}

\section{Intuition: Parametric Pruning}

\begin{intuitionbox}
\textbf{The Idea:} We don't know the optimal radius $r^*$, but we can guess it. For a guess $r$, build graph $G_r$ with edges of length $\leq r$. Can we cover all vertices with $k$ centers of radius $r$? This means finding a dominating set of size $k$ in $G_r$. Checking this is NP-hard, but we can approximate: find a maximal independent set in $G_r^2$ (square graph). If its size $\leq k$, its neighborhoods in $G_r$ cover everything.

\textbf{Why Square Graph?} If $M$ is an MIS in $G^2$, then every vertex is within distance 2 of $M$ in $G$. In metric spaces, distance 2 means radius $\leq 2r$ in the original metric. So we get a 2-approximation.
\end{intuitionbox}

\section{Problem Definition}

\begin{definition}[Metric $k$-Center]
Given complete graph with metric distances, choose $k$ centers minimizing the maximum distance from any vertex to its nearest center:
\[
\min_{S \subseteq V, |S|=k} \max_{v \in V} \min_{u \in S} \text{dist}(u, v)
\]
\end{definition}

This is the \emph{warehouse location} problem: minimize worst-case delivery distance.

\subsection{Parametric Pruning Algorithm}

\begin{algorithmdef}[Metric $k$-Center — Factor 2]
\begin{enumerate}
    \item Sort distinct edge weights; for each threshold $r$, build $G_r$ (edges $\leq r$)
    \item For each $G_r$, compute $G_r^2$ (square: edges for distance $\leq 2$)
    \item Find maximal independent set $M$ in $G_r^2$
    \item If $|M| \leq k$ and $M$'s neighborhoods in $G_r$ cover all vertices, output neighborhoods as centers with radius $2r$
\end{enumerate}
\end{algorithmdef}

\begin{theorem}
This algorithm achieves a factor-2 approximation for Metric $k$-Center.
\end{theorem}

\subsection{Weighted $k$-Center}

Vertices have weights; total weight of centers $\leq W$. Same parametric pruning with weight-aware selection gives factor-3 approximation.

\subsection{Other Approaches}

\begin{table}[H]
\centering
\caption{k-Center Algorithms}
\begin{tabular}{lccc}
\toprule
Algorithm & Approximation & Technique \\
\midrule
Parametric pruning (this chapter) & 2 & Combinatorial \\
Farthest-first traversal & 2 & Greedy \\
LP rounding & 2 & LP relaxation \\
High-dimensional $\ell_p$ & $1+\epsilon$ & Coresets \\
\bottomrule
\end{tabular}
\end{table}

\textbf{Why not exact?} $k$-Center is NP-hard even for $k=2$ (reduction from Dominating Set). Without metric assumption, no $\alpha(n)$-approximation exists.

\section{Python Implementation}

\begin{lstlisting}[style=python, caption=k-Center Parametric Pruning]
def build_graphs_by_threshold(graph):
    """Sorted set of edge weights and the subgraphs G_i of edges <= t."""
    edges = [w for u in graph for v, w in graph[u].items() if u < v]
    thresholds = sorted(set(edges))
    graphs = []
    for t in thresholds:
        g = {u: {} for u in graph}
        for u in graph:
            for v, w in graph[u].items():
                if w <= t:
                    g[u][v] = w
        graphs.append(g)
    return thresholds, graphs


def graph_square(graph):
    """Return G^2: edges between vertices reachable in <= 2 hops."""
    g2 = {u: {} for u in graph}
    for u in graph:
        g2[u] = dict(graph[u])                 # distance 1 neighbors
        for v in graph[u]:                     # distance 2 neighbors
            for w in graph[v]:
                if w != u:
                    s = graph[u][v] + graph[v][w]
                    if s < g2[u].get(w, float('inf')):
                        g2[u][w] = s
    return g2


def maximal_independent_set(graph):
    """Greedy maximal independent set."""
    indep, remaining = set(), set(graph)
    while remaining:
        v = next(iter(remaining))
        indep.add(v)
        remaining -= {v} | set(graph[v])
    return indep


def kcenter_parametric_pruning(graph, k):
    """2-approx for metric k-Center via parametric pruning (Algo 5.3)."""
    thresholds, graphs = build_graphs_by_threshold(graph)

    for t, g in zip(thresholds, graphs):
        mis = maximal_independent_set(graph_square(g))
        if len(mis) <= k:
            centers = set(mis)
            if len(centers) < k:
                remaining = set(graph) - centers
                centers.update(list(remaining)[:k - len(centers)])
            return centers, t

    return set(graph), max(thresholds)
\end{lstlisting}

\section{Applications}
\begin{itemize}
    \item \textbf{Facility location:} Fire stations, warehouses, cell towers
    \item \textbf{Clustering:} Minimize cluster radius ($k$-center = bottleneck $k$-means)
    \item \textbf{Load balancing:} Distribute tasks to minimize max load
    \item \textbf{Robotics:} Multi-robot coverage of an area
\end{itemize}

\subsection*{Real Example: Cell Tower Placement for 5G}
\textbf{Scenario:} Telecom placing $k=15$ new 5G towers in a city of $n=2,000$ potential sites. Each site covers a radius $r$. Goal: minimize $r$ so every subscriber is covered.
\textbf{Model:} Complete graph with metric distances (lat/long $\to$ haversine). $k$-center finds $k$ centers minimizing maximum distance to nearest center = maximum subscriber-tower distance.
\textbf{Why 2-approx matters:} If OPT radius = 800m, our towers have radius $\leq$ 1.6km. Guarantee ensures no subscriber is more than 2$\times$ the minimum possible distance from a tower.
\textbf{Real constraint:} Some sites are unavailable (zoning, cost). Remove those vertices from graph — parametric pruning handles this naturally.
\textbf{Comparison:} $k$-means minimizes \emph{sum} of squared distances (average case), $k$-center minimizes \emph{maximum} distance (worst case). For coverage guarantees, $k$-center is correct model.
\textbf{Weighted $k$-center (Ch. 5):} Towers have different capacities (weights). Total capacity budget $W$. 3-approx algorithm extends directly.

\section{Tight Example: Wheel Graph}
Center connected to $n$ outer vertices (weight 1), outer cycle weight 2. For $k=1$, optimal center = hub (radius 1). Algorithm may pick outer vertex (radius 2). Ratio = 2.

\section{Exercises}
\begin{exercise}
Show that without triangle inequality, $k$-center has no $\alpha(n)$-approximation.
\end{exercise}
\begin{exercise}
Modify Algorithm 5.3 to use minimal dominating set instead of MIS in $G_i^2$.
\end{exercise}

